% LSPV-1 REPORT::CUTFLOW85::2026-10-23T15:55:17-07:00
\documentclass[11pt]{article}
\usepackage[margin=.88in]{geometry}
\usepackage{amsmath,amssymb,amsthm,booktabs,url}
\setlength{\emergencystretch}{3em}
\newtheorem{theorem}{Theorem}
\title{CUTFLOW85: A Parameterized Two-Way Source-Transport Bound for Proper 3-SAT}
\author{LSPV-1: Universal Decider Cost Bridge}
\date{23 October 2026, Los Angeles}
\begin{document}\maketitle
\begin{center}
\textbf{P versus NP remains OPEN.}\\
\textbf{PROVED/EXACT only for the declared two-way local-head source model.}\\
Independent mathematical review pending.
\end{center}
\section{Target; ideas proposed before literature}
The target is a possible parameterized acquisition/compression/continuation bound over complete adaptive SAT policies, without assuming its premise. Three falsifiable mechanisms precede source review:
\begin{description}
\item[M1.] Remove the fixed-full-pass restriction of PASSMEM84 and prove a crossing-information bound for arbitrary two-way head reversals and arbitrary source-dependent symbolic method switches, charging every portable state bit.
\item[M2.] Build a source-validated batch comparator and an observed-evidence-triggered policy switch, testing matched SAT and mutated UNSAT sources.
\item[M3.] Falsify any universal transport claim through an interleaved encoding of exactly the same Boolean constraints and through paid random-access word addressing.
\end{description}
\noindent\textbf{CORPUS REVIEWED.} Root \path{main.tex}, AGENTS.md, CODEX.md, pnp-research skill, task/role tree, LSPV-1 PASSMEM84 and ADAPTIVEPROOFRACE83, MAIN85 ROOT-LABEL-QUOTIENT, MAIN108 WEIGHTED-TRANSCRIPT-COVER, Role 2 QUIET-EQUIV67 and SCC-ISOLATION68, Role 3 ADAPTIVE-COVER75, Role 4 SCC-AFFINE-LOCK70, and the earlier SEQUENTIAL-TAPE-001. The September 30 source already describes generic palindrome/equality crossing-sequence sums; PASSMEM84 already proves $(2p-1)s\ge k$ for multipass $E_k$. The precise new local source contribution is an arbitrary-turn parameterized source-cut theorem. No original general communication complexity theorem is claimed.

\noindent\textbf{LITERATURE CHECK.} Hennie established the local crossing-sequence method; Kushilevitz--Nisan give equality fooling-set/rectangle arguments; Ibarra--Moran study restricted tape/offline time-space models. Aspvall--Plass--Tarjan's SCC 2SAT and Darmann--D\"ocker's strict signed-monotone exact-(2,2) endpoint are hostile controls. None supplies an all-RAM SAT lower bound.

\section{Exact source and shared-input Circuit-SAT construction}
Let $\lambda_t(x)=x$ for $t=1$ and $\lambda_t(x)=\neg x$ for $t=0$, and put
\[
Q(\ell;a,b)=\bigwedge_{\epsilon,\delta\in\{-1,+1\}}
 (\ell\vee a^\epsilon\vee b^\delta).
\]
Take $u,v\in\{0,1\}^k$ and five distinct source variables $(x_i,a_i,b_i,c_i,d_i)$ per coordinate. Define
\[
E_k(u,v)=
\left(\bigwedge_{i=1}^k Q(\lambda_{u_i}(x_i);a_i,b_i)\right)
\land
\left(\bigwedge_{i=1}^k Q(\lambda_{v_i}(x_i);c_i,d_i)\right).
\]
All first-half $u$ quartets precede the second-half $v$ quartets on the physical source-word tape.
Since $Q(\ell;a,b)\equiv\ell$, the source is proper unique three-literal CNF with $n=5k$ and $m=8k$, and
\[
\boxed{E_k(u,v)\in\mathsf{SAT}\iff u=v},\qquad
\#SAT(E_k(u,u))=2^{4k}.
\]
Under fixed-width encoding with two public header words and four words per clause, word width
$w=\lceil\log_2(64(13k+1))\rceil$, total word cells $W=32k+2$, and encoded bits $N=Ww$.
Only two groups of $k$ signs vary. The input is NOT proper signed exact-(2,2), since each $x_i$ appears eight times. Strict signed exact-(2,2) $H_g$ is retained as a separate control.
A shared-named-input NOT1/OR2/AND2 evaluation circuit has $G=n+3m-1=29k-1$ gates and
$J=n+6m-2=53k-2$ directed gate-input incidences. Functional gate-definition CNF has $n+G=34k-1$ variables and $2n+9m-2=82k-2$ clauses including output TRUE. Each auxiliary gate value is fixed by the original source assignment. The backward AND-demand, OR-choice selector and $K_{2,2}$ occurrence conflict graph are exact reformulations when identities are preserved, not mandatory work or SAT-solving circuit lower bounds.

\section{Model, independently defined parameter, and proof}
Fix a uniform deterministic exact algorithm $A$ on every canonical source. Its read-only input head moves at most one \emph{source-word cell} per charged step, left, right or stationary. Every cache, partial assignment, learned clause, work head location, rule selector, timer, alternate representation and bit of persistent computation outside the source must reside in a portable configuration with at most
$M_A(k)$ possibilities on length-$k$ canonical sources. Write
$\sigma_A(k)=\max(1,\lceil\log_2 M_A(k)\rceil)$, which is at most charged work bits $s_A(k)$ plus an algorithm-dependent fixed-control constant $c_A$. There are no free random seeks, source writes, outside memory or uncharged clocks. The internal policy may switch arbitrarily often among all computable symbolic rules within these exact source channels.

For every public cut after the $j$th first-half quartet and after the $j$th second-half quartet, define the structural cross-side fooling dimensions
\[
h_{u,j}=j,\qquad h_{v,j}=k-j,\qquad
\mathcal H(E_k)=\sum_{j=1}^{k-1}(h_{u,j}+h_{v,j})=k(k-1).
\]
They are defined by the source layout and its exact SAT semantics independently of measured runtime or actual internal solver choices.

\begin{theorem}[CUTFLOW85: deterministic adaptive two-way transport, PROVED/EXACT]
For every such complete exact algorithm $A$, every $k\ge3$, and $X$ uniform in $\{0,1\}^k$,
\[
\boxed{\mathbb E_X[T_A(E_k(X,X))]
\ge \frac{(k-1)(k-2)}{\sigma_A(k)+1}.}
\]
Thus some algorithm-dependent satisfying source attains this bound. Every positive canonical source also forces the head to inspect a sign in the last $v$ quartet, at word index at least $32k-13$, so for the relevant worst-case time and complete live memory:
\[
\boxed{T_A^{\max}(k)=
\Omega\!\left(k+\frac{k^2}{s_A(k)+c_A+1}\right).}
\]
The constants in $\Omega$ depend only on the fixed machine and source-word convention.
\end{theorem}
\begin{proof}
For the cut after $u_j$, fix the suffix $X_{j+1..k}=t$. The $2^j$ accepted diagonal sources
$E_k((a,t),(a,t))$, $a\in\{0,1\}^j$, have distinct \emph{complete crossing-configuration transcripts}:
if the transcripts for $a\ne b$ agreed, deterministic read-only cut-splicing would accept the mixed
$E_k((a,t),(b,t))$, an UNSAT instance. The transcript can be encoded prefix-free by a bit 1 and the full
$\sigma_A$-bit portable configuration for each crossing, followed by a 0 endmarker. If $C_{u,j}$ is the
number of crossings, code length equals $(\sigma_A+1)C_{u,j}+1$. Prefix-free source coding applied to
uniform $a$ yields $(\sigma_A+1)\mathbb E[C_{u,j}\mid t]+1\ge j$, hence
$\mathbb E C_{u,j}\ge(j-1)/(\sigma_A+1)$ after averaging over $t$.
For the cut after $v_j$, fix the prefix $X_{1..j}$ instead. Splicing distinct accepted suffixes
must similarly produce a rejected mixed input, giving
$\mathbb E C_{v,j}\ge(k-j-1)/(\sigma_A+1)$.
All $2(k-1)$ physical quartet-boundary cuts are distinct. A unit local input-head move crosses
at most one of them. Therefore
\[
\mathbb E T_A\ge\sum_{j=1}^{k-1}
\frac{(j-1)+(k-j-1)}{\sigma_A+1}
=\frac{(k-1)(k-2)}{\sigma_A+1}.
\]
Flipping the last $v$ bit alone changes SAT to UNSAT and changes only first-literal sign words in the
last $v$ quartet. On every positive source an exact decider must inspect at least one of them, whose
earliest position is word index $32k-13$ when the head starts at zero. This yields the independent
linear source-travel bound and the stated joint result. Internal representations and adaptive switches
are part of the crossing configurations, rather than being prescribed by the proof.
\end{proof}

\section{Charged constructive match and two decisive escapes}
A complete canonical source-grammar recognizer makes one paid sequential word sweep,
with $O(\log k)$ bits of indices. A deterministic batch policy retains $B$ consecutive $u$ signs,
moves its physical source head to the corresponding $v$ signs, compares, and repeats.
It returns UNSAT immediately upon an observed mismatch and SAT only after every coordinate agrees.
Its exact restricted-family upper bound is
\[
T=O(k+k^2/B),\qquad s\le B+O(\log k)\ \text{bits}.
\]
An adaptive-first variant checks one pair with one saved sign, and switches to a larger $B$ after
observed agreement. The true label is never given in advance. Failed recognizers on arbitrary CNF
require a separately charged complete fallback, and no all-input small-space promise is made.

\noindent\textbf{FALSIFIED/RETIRED universal transport.} A random-access word-RAM validates the public source grammar and compares addressed signs in $O(k)$ word probes with $O(\log k)$ space. Alternatively, the \emph{same eight source clauses per coordinate} in the interleaved order $Q(u_1),Q(v_1),\ldots,Q(u_k),Q(v_k)$ admit local-head one-sweep comparison in $O(k)$ motion and $O(\log k)$ bits. The Boolean formula, variable identities, model count and gate counts are unchanged, but the physical cut profile collapses. Reordering a \emph{given} grouped tape is not free and must be charged; an originally interleaved source is a separate legal encoding. Thus the cut-entropy parameter is access-layout relative, not a semantic SAT-hardness invariant.

\section{Actual adaptive state ledger and finite audit}
\begin{description}
\item[Before:] Source records exist on a read-only local word tape; only public metadata and an allowed payload bound are available. A paid grammar scan learns $k$, validates quartet identities and locates relevant sign records without retaining the source.
\item[Action:] Compare fixed-$B$ batching against the adaptive first-pair probe, with interleaved input and addressed-word access as model-mismatched hostile controls. Selection, head travel, reads, policy evaluation, failed branches and switches are charged.
\item[After:] Retain at most $B$ pending signs plus logarithmic indices. An observed mismatch returns UNSAT; after all comparisons match, SAT follows from the exact quartet equivalence. All source reconstruction if a witness is requested must be separately replayed and verified; the tested service only decides.
\item[Peak:] Explicit payload $B$ and $O(\log k)$ full-control bits for these algorithms; the theorem measures the \emph{entire} portable configuration $\sigma_A$, not merely payload.
\end{description}
\noindent\textbf{COMPUTATIONAL.} Reproducible standard-library Python with seed 202610231555 passed 353 assertions, zero failures: 20 exhaustive small $(u,v)$ source assignments and SAT labels; 192 grouped SAT/UNSAT fixed/adaptive local-head executions across $k=8,16,32,64$, four label/mismatch positions and six budgets; three strict signed exact-(2,2) $H_g$ controls; malformed source rejection; direct-access and interleaved comparator controls; and exact $k(k-1)$ cut-profile arithmetic.
For $k=64$, $B=8$, grouped SAT cost was 22,417 movement-plus-read events, and adaptive-first SAT 24,433. An early UNSAT mutation took 7,178 versus 7,171 events, with the adaptive-first policy retaining one rather than eight signs before rejection. Grouped SAT with $B=1$ took 135,313; with $B=64$ took 8,305. Interleaving required 8,305 with one retained sign. Direct addressed access used 2,178 read probes including syntax validation. These are instrumented abstract costs, not physical CPU instruction counts or asymptotic hardness evidence.
Raw results are in \path{research/agents/lspv-1/2026-10-23_1555_CUTFLOW85_results_summary.json}; the compact verifier is in the same directory.

\section{No-double-counting, peer integration, and exact next gap}
Cut-crossing counts are subsets of real moves, so summing distinct cuts is legitimate. Adding imagined clause-propagation costs to that theorem without disjointness would not be. The parameter is a spatial information profile across source records and may be inspected analytically even when the solver never constructs it. This is a restricted time--space lower bound, not an exponential one. Neither $P=NP$ nor $P\ne NP$ follows.

\noindent\textbf{WHAT CHANGED FROM ROOT MAIN.} PASSMEM84's $p$-pass lower bound is replaced by an arbitrary-turn two-way theorem for the identical source, preserving all adaptive internal symbolic methods subject to the declared fully charged local head and live state. This instantiates rather than independently discovers the older SEQUENTIAL-TAPE-001 crossing-battery concept. Root MAIN85's two-state SAT quotient remains decisive against an all-algorithm interpretation. Role 2 SCC collapses, Role 3 cover, and Role 4 SCC-affine switches remain genuine constructive hostile controls.

\noindent\textbf{Peer handoff.} LSPV-2: compare grouped/interleaved backward TRUE-gate source representations. LSPV-3: red-team crossing splicing against hidden control, real-time clocks, random access, multiple heads and advice. LSPV-4: SAT/UNSAT method-switching experiments and source grammar mutation. LSPV-5: retain the $\Omega(k^2/s)$ claim only in the specified one-head local-read model; note that it is polynomial and source-order sensitive. Preserve all five active schedules unchanged.

\noindent\textbf{OPEN first missing implication.} Find a noncircular semantic or source-grounded parameterized obstruction that survives arbitrary \emph{random-access} adaptive SAT algorithms including SCC, affine/Horn, DPLL, learning, caching and direct decision. Prove separately its fully charged work inequality and its parameter growth on arbitrarily large algorithm-dependent inputs. The required universal quantifiers are
\[
\forall A\ \forall c>0\ \forall N_0\ \exists F_A:
|F_A|\ge N_0\ \land\ T_A(F_A)>|F_A|^c.
\]
Both the all-decider bridge and this superpolynomial conclusion remain OPEN.

\begin{thebibliography}{5}
\bibitem{hennie} F. C. Hennie, One-tape, off-line Turing machine computations, \emph{Information and Control} 8 (1965), 553--578.
\bibitem{kn} E. Kushilevitz and N. Nisan, \emph{Communication Complexity}, Cambridge University Press, 1997.
\bibitem{im} O. H. Ibarra and S. Moran, Some Time-Space Tradeoff Results Concerning Single-Tape and Offline TM's, \emph{SIAM Journal on Computing}, 1983.
\bibitem{apt} B. Aspvall, M. F. Plass and R. E. Tarjan, \emph{Information Processing Letters} 8 (1979), 121--123.
\bibitem{dock} A. Darmann and J. D\"ocker, \emph{Discrete Applied Mathematics} 292 (2021), 45--58.
\end{thebibliography}
\end{document}
% END_OF_VERIFIED_REPORT::LSPV-1::CUTFLOW85